\documentclass[12pt,a4paper]{article}
\usepackage[utf8]{inputenc}
\usepackage[italian]{babel}
\usepackage{geometry}
\usepackage{graphicx}
\usepackage{titlesec}
\usepackage{tabularx}
\usepackage{enumitem}

\geometry{margin=2.5cm}

\begin{document}

% --- PAGINA 1: COPERTINA ---
\begin{titlepage}
    \centering
    \vspace*{2cm}
    {\Huge \textbf{Nome progetto} \par}
    \vspace{1cm}
    {\Large Deliverable D1 \par}
    \vspace{3cm}
    
    {\large \textbf{Gruppo n°:} 20 \par}
    \vspace{1.5cm}
    
    {\large \textbf{Componenti:} \par}
    \vspace{0.5cm}
    \begin{tabular}{rl}
        \textbf{Studente:} & Binco Francesco (Matricola: ) \\
        \textbf{Studente:} & Cortese Riccardo (Matricola: 244656) \\
        \textbf{Studente:} & Giovagnetti Lorenzo (Matricola: ) \\
    \end{tabular}
    
    \vfill
    {\large Progetto di Ingegneria del Software - Prof. Sandro Luigi Fiore \par}
    \vspace{1.5cm}
    {\large Anno Accademico 2025/2026 \par}
    \vspace{2cm}
\end{titlepage}

% --- PAGINA 2: INDICE ---
\newpage
\pagenumbering{roman} % Numerazione romana per l'indice
\tableofcontents
\newpage

% --- PAGINA 3 E SUCCESSIVE: CONTENUTO ---
\pagenumbering{arabic} % Torna alla numerazione araba (1, 2, 3...)
\setcounter{page}{1}   % Ricomincia il conteggio da 1

\section{Obiettivi del progetto}
L'obiettivo principale del sistema è lo sviluppo di un'applicazione che migliori l'efficienza della raccolta differenziata, facilitare il corretto smaltimento dei rifiuti da parte dei cittadini e ottimizare il lavoro degli operatori ecologici.

\section{Attori del sistema}
% Nota: Inserire qui una mindmap (consigliata) [cite: 10]
\subsection{Utenti}
\begin{itemize}
    \item \textbf{Utente non autenticato:} utente base con funzionalità limitate[cite: 12].
\end{itemize}

\subsection{Sistemi esterni}
[Indicare eventuali sistemi esterni] [cite: 13]

\section{Requisiti funzionali per attore}
% Nota: Almeno 20 totali [cite: 15]
\subsection{Utente anonimo o non autenticato}
\begin{itemize}
    \item \textbf{RF1:} L'applicazione deve fornire agli utenti anonimi una breve spiegazione riguardo le funzionalità dell'applicazione e del progetto[cite: 18].
\end{itemize}

\section{Requisiti non funzionali}
% Nota: Almeno 10 totali [cite: 20]
\begin{itemize}
    \item \textbf{RNF1:} La velocità di risposta delle funzioni di ricerca del sistema deve essere inferiore ai 5 secondi[cite: 21].
\end{itemize}

\section{Use case diagrams e casi d'uso}
% Nota: Rappresentare diagramma e almeno 5 descrizioni [cite: 23, 24]
\begin{table}[h]
\centering
\begin{tabularx}{\textwidth}{|l|X|}
\hline
\textbf{ID e Nome} & UC1: [Nome Caso d'Uso] \\ \hline
\textbf{Attori} & \dots \\ \hline
\textbf{Descrizione} & [Descrizione dettagliata con tabella] [cite: 25] \\ \hline
\end{tabularx}
\end{table}

\section{Diagramma BPMN}
% Nota: Almeno un diagramma BPMN (as-is, to-be o dettaglio) [cite: 27, 28]
\textit{// Inserire qui il diagramma BPMN.}

\end{document}